\section{Frekvens og stabilitet}

\subsection{Bode-plot og P-gain / Bode, gain margin, phase margin}
\label{sec:bode}

\paragraph{Brug når}
Opgaven giver en transferfunktion eller Bode-plot og spørger om poler/zeros,
stabilitet, \(K_P\), gain margin eller phase margin.

\paragraph{Input fra opgaven}
Break frequencies, magnitude i dB, fasekurve, controllergain og ønsket margin.

\paragraph{Metode}
\begin{enumerate}
\item Marker startgain, integratorer og hver ændring i magnitudeslope.
\item Brug faseændringen til at skelne LHP fra RHP-singulariteter.
\item Find \(\omega_c\) ved 0 dB eller \(\omega_\pi\) ved \(-180^\circ\).
\item Beregn margin eller gainkorrektion; konvertér dB til lineær gain.
\item Kontrollér stabilitetskonklusionen med plantens open-loop-stabilitet.
\end{enumerate}

\paragraph{Formler}
\[
M_{\mathrm{dB}}=20\log_{10}|G(\jw)|,\qquad
|G|=10^{M_{\mathrm{dB}}/20},
\]
\[
|L(\mathrm{j}\omega_c)|=1,\qquad
\gamma_M=180^\circ+\angle L(\mathrm{j}\omega_c),\qquad
K_M=\frac1{|L(\mathrm{j}\omega_\pi)|}.
\]
\begin{center}
\begin{tabular}{@{}lcc@{}}
\toprule
Singularitet & Magnitudeslope efter break & Samlet phasebidrag\\
\midrule
LHP-pol & \(-20\) dB/dec & \(-90^\circ\)\\
LHP-zero & \(+20\) dB/dec & \(+90^\circ\)\\
RHP-pol & \(-20\) dB/dec & \(+90^\circ\)\\
RHP-zero & \(+20\) dB/dec & \(-90^\circ\)\\
\bottomrule
\end{tabular}
\end{center}

\paragraph{Symboler}
\(\omega_c\) er gain crossover, \(\omega_\pi\) phase crossover,
\(\gamma_M\) phase margin og \(K_M\) gain margin.

\paragraph{Antagelser}
Almindelig ``positiv margin giver stabilt lukket loop'' anvendes direkte for de
open-loop-stabile/minimum-phase situationer, som kriteriet forudsætter; RHP-planter
behandles med Nyquist nedenfor.

\paragraph{Hurtigtjek}
Positiv \(K_P\) løfter magnitude med \(20\log_{10}K_P\), men ændrer ikke fase.
Netto højfrekvensslope skal passe med relativ grad.

\paragraph{Typiske fælder}
At bruge 38.9 dB som gain 38.9 i stedet for \(10^{38.9/20}\approx88\) (F21 Q6);
at afgøre RHP/LHP uden phaseplot.

\paragraph{Python}
\code{evaluate\_transfer\_function(num,den,omega)} kan evaluere en kendt model;
\code{closed\_loop\_poles(num,den,K)} kan kontrollere foreslået P-gain. Brugeren skal
selv aflæse plots og verificere loopstruktur.

\subsection{Nyquist-stabilitet for ustabil plant / encirclement, RHP poles}
\label{sec:nyquist}

\paragraph{Brug når}
Plantens open-loop har RHP-poler, eller opgaven angiver Nyquistkurvens skæring
og retning omkring \((-1,0)\).

\paragraph{Input fra opgaven}
Antal RHP-poler \(P\), netto kurveretning og realakseskæring samt controllergainens fortegn.

\paragraph{Metode}
\begin{enumerate}
\item Find \(P\), antal open-loop RHP-poler.
\item Tæl clockwise encirclements \(N\) af \(-1\) med fortegn fra figuren.
\item Brug \(Z=P+N\); stabilt lukket loop kræver \(Z=0\).
\item Skalér Nyquistkurven med positiv gain eller vurder spejling ved negativ gain.
\item Kontrollér at stabiliseringsgain ikke blot rammer marginalpunktet.
\end{enumerate}

\paragraph{Formler}
\[
Z=P+N,\qquad Z=0\ \text{for closed-loop-stabilitet.}
\]
Har \(P=1\), kræves ét counter-clockwise encirclement, dvs. \(N=-1\) i
forelæsningernes konvention. Hvis den relevante CCW-skæring er \(-a\), kræves
typisk \(K_P>1/a\) for positiv radial skalering.

\paragraph{Symboler}
\(P\) er open-loop RHP-poler; \(N\) netto clockwise encirclements; \(Z\) closed-loop
RHP-poler.

\paragraph{Antagelser}
Nyquistkonturen og eventuelle imaginæraksepoler håndteres som i den givne figur;
den angivne skæring er den stabiliserende passage.

\paragraph{Hurtigtjek}
F21 Q13: \(a=0.0111\), så \(1/a\approx90.1\); valget \(K_P=100\) passerer \(-1\).

\paragraph{Typiske fælder}
At vælge lille positiv gain for en open-loop-ustabil plant; at vende CW/CCW-fortegn.

\paragraph{Python}
Ingen funktion automatiserer encirclement fra figuren. Efter manuel stabiliseringsbeslutning
kan \code{closed\_loop\_poles} bruges, hvis en fuld analytisk \(G(s)\) er angivet.

\subsection{Phase margin fra Nyquist-punkt / unit-circle point}

\paragraph{Brug når}
Nyquistkurvens punkt på enhedscirklen ved crossover opgives numerisk.

\paragraph{Input fra opgaven}
Punktet \((x,y)=(\Re L,\Im L)\), der svarer til \(|L|=1\).

\paragraph{Metode}
\begin{enumerate}
\item Beregn vinklen med \(\operatorname{atan2}(y,x)\), ikke enkel arctan uden kvadrant.
\item Beregn \(\gamma_M=180^\circ+\phi_c\).
\item Sammenlign størrelsesorden og svarenhed med valgmulighederne.
\end{enumerate}

\paragraph{Formler}
\[
\phi_c=\operatorname{atan2}(y,x),\qquad \gamma_M=180^\circ+\phi_c .
\]

\paragraph{Symboler}
\(x,y\) er real- og imaginærdel af open-loop responsen; \(\phi_c\) er fase ved crossover.

\paragraph{Antagelser}
Punktet er korrekt crossoverpassage og vinklen udtrykkes i grader.

\paragraph{Hurtigtjek}
F21 Q14: \((0.134,-0.99)\) ligger i fjerde kvadrant, så \(\phi_c\) er omkring
\(-82^\circ\) og marginen omkring \(98^\circ\).

\paragraph{Typiske fælder}
At aflevere \(\phi_c\), radianværdien eller \(-\phi_c\) som margin.

\paragraph{Python}
\code{phase\_margin\_from\_point(real\_part,imaginary\_part)} returnerer marginen,
når du har aflæst det rette punkt.
